\section{Joined source estimates and a uniform parameter envelope}
\label{sec:parameter-envelope}

We fix the dimension throughout this section. The Sobolev base order
$q=2d+5$ is fixed before any stage or truncation is chosen. Increasing this
order changes constants and polynomial degrees only.

\subsection{The variational history inverse}

For every forced high grade, the inverse consists of a Dirichlet problem
before activation and a forward evolution after activation. The Dirichlet
coercivity estimate gives polynomial dependence on the parent coefficients
on the history interval. On the forward interval this dependence follows
from the relative propagator estimate. Compact-interval existence of a
forward fundamental matrix alone does not provide either quantitative bound.

The quantitative estimate is established in the transverse space
$U=m_0^\perp$, of dimension $d-1$. Let
\[
 Q(t,x)=F(t,x)|_U,\qquad Q_t=Q_1,\qquad
 Q_{tt}=-\mathcal H Q,\qquad
 \|Qv\|^2\geq c\|v\|^2,\quad c>0.
\]
The symmetric matrix $\mathcal H$ is the actual pressure Hessian along
the old particle map. For a positive activation time $0<t_0\leq T$, suppose on $[0,t_0]$
\[
 \langle\mathcal H v,v\rangle\leq K\|v\|^2,\qquad
 0\leq K,\qquad Kt_0^2/2\leq1/2.
\]
These are the inherited one-sided pressure bounds. All norms below can
also be taken in spatial-angular $L^2$; the multiplication coefficients
have uniform pointwise bounds.

For a forcing $f$, solve for $z\in H^1_0((0,t_0);U)$:
\begin{equation}\label{eq:joined-history-form}
 \int_0^{t_0}\!
 \big\langle(Qz)',(Q\zeta)'\big\rangle
 -\big\langle\mathcal H Qz,Q\zeta\big\rangle\,dt
 =-\int_0^{t_0}\langle f,Q\zeta\rangle\,dt
 \quad(\zeta\in H^1_0((0,t_0);U)).
\end{equation}
Its physical Jacobi displacement is $y=Qz$; its physical high velocity is
\begin{equation}\label{eq:joined-history-high}
 b=Qz',
\end{equation}
which is generally different from $y'$. In particular, $z(0)=z(t_0)=0$
does not assert that $b(0)$ or $b(t_0)$ vanishes.

\begin{lemma}[Polynomial history inverse]\label{lem:joined-history-polynomial}
The solution of \eqref{eq:joined-history-form}, its high velocity,
its time derivative and its terminal high velocity have bounds
polynomial in
\[
 1+t_0+t_0^{-1}+c^{-1}+
 \|Q\|_\infty+\|Q_1\|_\infty+\|\mathcal H\|_\infty .
\]
Their fixed-order spatial-angular Sobolev bounds are also polynomial in
the corresponding fixed-order coefficient bounds. The assertion holds
in every finite transverse dimension.
\end{lemma}
\begin{proof}
Put $C_0=\|Q\|_\infty$, $C_1=\|Q_1\|_\infty$ and
$L=(Q^*Q)^{-1}Q^*$. Then $LQ=I$ and
\[
 \|L\|\leq c^{-1}C_0,\qquad
 \|L'\|\leq c^{-1}C_1+2c^{-2}C_0^2C_1.
\]
Consequently, with
\[
 A=1+c^{-1}C_0+
       t_0(c^{-1}C_1+2c^{-2}C_0^2C_1),
\]
the identity $z=Ly$ and the elementary inequality
$\|y\|_{L^2}\leq t_0\|y'\|_{L^2}$ give
$\|z'\|_{L^2}\leq A\|y'\|_{L^2}$.
The sharper primitive bound
$\|y\|_{L^2}^2\leq(t_0^2/2)\|y'\|_{L^2}^2$ gives
\[
 \int_0^{t_0}(|y'|^2-\langle\mathcal Hy,y\rangle)
 \geq\tfrac12\|y'\|_{L^2}^2
 \geq(2A^2)^{-1}\|z'\|_{L^2}^2.
\]
Thus the coercive theorem on the fixed Hilbert space $H^1_0$ constructs
the inverse. Testing with $z$ also gives
\[
 \|y'\|_{L^2}\leq2t_0\|f\|_{L^2},\qquad
 \|z'\|_{L^2}\leq2At_0\|f\|_{L^2}.
\]
Both bounds are polynomial in $C_0$, $C_1$ and $\|\mathcal H\|$.

Integration by parts and $Q''=-\mathcal HQ$ give
\[
 Q^*(Qz''+2Q_1z')=Q^*f.
\]
Therefore the actual acceleration is
\begin{equation}\label{eq:joined-gram-acceleration}
 z''=(Q^*Q)^{-1}\{Q^*f-2Q^*Q_1z'\},
\end{equation}
and
\[
 \|z''\|_{L^2}\leq
 c^{-1}C_0(\|f\|_{L^2}+2C_1\|z'\|_{L^2}).
\]
The time trace inequality
\[
 \|z'\|_{C_t}\leq
 t_0^{-1/2}\|z'\|_{L^2}+t_0^{1/2}\|z''\|_{L^2}
\]
bounds both endpoints polynomially. For continuous $f$,
\eqref{eq:joined-gram-acceleration} gives continuous $z''$, and
$b'=Q_1z'+Qz''$ is the actual derivative on the closed interval.
In particular the forward initial datum $b(t_0)$ is a bounded trace
of the solution of \eqref{eq:joined-history-form}.

Spatial derivatives are obtained by differentiating the operator on
the fixed space $H^1_0((0,t_0);U)$. Every commutator contains only
coefficient derivatives and a lower derivative of the same solution.
Applying the same coercive inverse proves the fixed-order estimates
by finite induction. For example, a sufficient inverse cost is
\[
 I_0=I,\qquad I_{r+1}=I+I_r+2^rB I_r^2,
\]
where $I$ bounds the base inverse and $B$ the coefficient block.
This is a polynomial after the fixed number $q$ of steps.
The Gram inverse in \eqref{eq:joined-gram-acceleration} has the same
positive-order commutator proof. These arguments use adjoints, norms
and Hilbert-space coercivity, and contain no three-dimensional operation.
\end{proof}

On $[t_0,T]$ solve the actual forward transverse equation with initial
datum $b(t_0)$ obtained above. Its value and derivative match those of
the history, by the same equation and forcing at the junction.
The algebraic normal pressure is reconstructed from this joined field.
The profile is $\gamma(t)=\alpha$ on the history interval and
$\gamma(t)=\alpha g(t)$ afterwards, where $g(t_0)=1$ and $\alpha>0$.
Normalization by $\gamma$ in spatial word estimates never differentiates
$\gamma$ in time.

\subsection{The full transverse relative estimate}
\label{subsec:full-transverse-relative}

We give the additional-dimensional calculation explicitly. At activation
let $p,q$ be the old unit normal and unit velocity, and put
\[
 a=\langle Bq,p\rangle>0,\qquad
 \beta=\frac{\|\Pi_{\{p,q\}^{\perp}}Bq\|}{a},\qquad
 \sigma=\sqrt\beta,\qquad \varepsilon=\sqrt{a/h}.
\]
Here $h$ is the old shear height. We use $0<\sigma\leq1/4$,
$1/2\leq a\leq2$ and $0<\varepsilon\leq1$.
Choose the third unit vector $n$ in the direction of
$\Pi_{\{p,q\}^{\perp}}Bq$; the positive tilt makes it nonzero.
Write $E=\operatorname{span}\{p,q,n\}$ and
$\mathcal E=E^\perp$. The center reflection fixes $E$ and reverses
$\mathcal E$. The actual center coefficient therefore preserves this
splitting. Neighboring coefficients need not preserve it.

Use the inherited bounds $G\geq1$, $\|B\|\leq G$ and
$\|B_t\|\leq G^2$. Transport the entire orthonormal frame by its skew
connection $S$. In this frame, for the old background matrix $B$,
\[
 S=p\otimes B^*p-B^*p\otimes p+
       q\otimes\Pi_{\{p,q\}^{\perp}}Bq
       -\Pi_{\{p,q\}^{\perp}}Bq\otimes q .
\]
Consequently $\|S\|\leq4G$ if $\|B\|\leq G$, and the exact column identity is
\begin{equation}\label{eq:full-connection-column}
 (B+S)q=2\langle Bq,p\rangle p+\langle Bq,q\rangle q.
\end{equation}
The normal-normal connection is zero. Differentiating the transported
background gives norm at most $G^2+2(4G)G=9G^2$.
At activation its column outside the $q$ diagonal is
$a(p+\beta n)$. At later physical time $t$ its defect is therefore at
most $9G^2|t-t_0|$. These statements concern the complete vector column,
including all additional components.

Use fast time $\tau=a(t-t_0)/\varepsilon$. Only the actual interval
$0\leq\tau\leq\tau_{\rm end}$ is needed; let
$\tau_{\rm end}\leq\Theta$, where $\Theta\geq1$ is the numerical horizon.
No actual coefficient is extended to a longer interval.
At a neighboring label the actual gradient in this same frame is
\[
 M=B+h(t)\,q\otimes p+\mathcal R,\qquad
 \|\mathcal R\|\leq e_{\rm old}+\rho L_M .
\]
The rank-one term uses the center pair; the neighboring discrepancy is
included in $\mathcal R$.

Define the diagonal maps
\[
 D_r=\operatorname{diag}(1,\varepsilon,1,I_{\mathcal E}),\qquad
 D_v=\operatorname{diag}(\varepsilon,1,\varepsilon,
                                      \varepsilon I_{\mathcal E}).
\]
They satisfy $D_rD_v=\varepsilon I$ and both have norm at most one.
The actual ray and high velocity, including every extra coordinate, are
\begin{equation}\label{eq:ambient-scaled-states}
 m=s_0 O(t)D_r y,\quad y=(P,Q,N,r),\qquad
 v=O(t)D_v z,\quad z=(U,V,W,Z),\quad r,Z\in\mathcal E.
\end{equation}
The primary uses one fixed coordinate endpoint at every label.
More precisely, if the endpoint lemma selects the physical vector
$Y_{\rm phys}\in m(t_0,0)^\perp$, set
\[
 Y_{\rm coord}=R_\perp^*F(t_0,0)^{-1}Y_{\rm phys}.
\]
At the center, $F(t_0,0)R_\perp Y_{\rm coord}=Y_{\rm phys}$.
At every other label the history has the same coordinate endpoint
$Y_{\rm coord}$; no endpoint is selected again. Its norm is at most
$\|F(t_0,0)^{-1}\|\|Y_{\rm phys}\|$. Thus $Y_*$ in the history
Lipschitz estimates below includes this old inverse-frame factor.

For $t_0>0$, the unlocalized history is solved pointwise in the finite-dimensional
transverse space and has uniform spatial derivative bounds. Its
cylinder $L^2$ realization is obtained after multiplication by the
compact seed. To verify this directly, put
$Z(y,\theta)=\chi_{\rm in}(y)f_\delta(\theta)Y_{\rm coord}$
and $z_{\rm lin}(t,y,\theta)=tZ(y,\theta)/t_0$.
The zero-trace remainder $w=z-z_{\rm lin}$ solves
\eqref{eq:joined-history-form} with forcing
\[
 f_{\rm lin}=-(2/t_0)Q_1Z.
\]
This identity follows from $Q_{tt}=-\mathcal HQ$ and
$z_{\rm lin,tt}=0$. The forcing is a continuous path in every
cylinder Sobolev space, is angular-mean zero, and has compact spatial
support. Thus the Hilbert-space inverse applies. Its uniqueness
identifies the result with the localized pointwise history. All
cylinder estimates concern this localized field. For $t_0=0$ the
prescribed forward primary is used instead of the history construction.

Here $s_0>0$ and $O(t)$ is the actual transported orthogonal frame.
The velocity $v$ is the unamplified selected high solution. The packet
uses $\alpha v$, with $\alpha>0$ chosen from its actual target size below.
Set
\[
 A=a^{-1}D_rMD_v,\qquad C=a^{-1}D_r(M+S)D_v .
\]
Direct substitution into the ray and high equations gives
\begin{equation}\label{eq:ambient-exact-generator}
 y'=(C^*-2A^*)y,\qquad
 z'=\mathscr B_\varepsilon(y)z,\qquad
 \mathscr B_\varepsilon(y)
  =-C+2\frac{D_r^2y}{|D_ry|^2}\otimes A^*y .
\end{equation}
For example, the pressure numerator transforms by
$\langle D_ry,MD_vz\rangle=a\langle y,Az\rangle$.
Thus \eqref{eq:ambient-exact-generator} is the actual rescaled equation;
no extra-mode equation has been suppressed.

Let
\[
 A_0=(p+\beta n)\otimes q+q\otimes p,\qquad
 C_0=2p\otimes q+q\otimes p .
\]
The only large shear scales to
$(\varepsilon^2h(t)/a)q\otimes p$.
For any matrix $L$, separating its $q$ input column gives
\[
 a^{-1}D_rLD_v
 =(\varepsilon/a)D_rL+
          ((1-\varepsilon)/a)(D_rLq)\otimes q.
\]
Together with \eqref{eq:full-connection-column}, this proves
\begin{equation}\label{eq:ambient-matrix-gap}
 \max(\|A-A_0\|,\|C-C_0\|)
 \leq \left|\frac{\varepsilon^2h(t)}a-1\right|
 +\frac{11\varepsilon G+18G^2|t-t_0|+
                    e_{\rm old}+\rho L_M}{a}.
\end{equation}
One obtains the slightly better coefficient 3 in place of 11 for $A$;
the displayed common bound suffices. The norm of an additional input
column has the small factor $\varepsilon$: the shear is zero on
$\{p,q\}^\perp$. In particular there is no estimate of an extra column
by the full large shear. Nor does matrix scaling introduce
$\varepsilon^{-1}$, since both diagonal maps have norm at most one.

The logarithmic derivative of $h(t)$ is bounded by $2G$.
For $G\varepsilon\Theta/a$ small, its contribution to
\eqref{eq:ambient-matrix-gap} is $O(G\varepsilon\Theta/a)$.
Choose a single error parameter $e\geq\varepsilon$ large enough that
the right side of \eqref{eq:ambient-matrix-gap} is at most $2e$,
and that the initial ray error is at most $e$.
It is bounded by a fixed constant times
\begin{equation}\label{eq:ambient-error-input}
 \varepsilon\Theta G^2+
 e_{\rm old}+
 \rho\left(L_M+\frac{3L_m}{s_0\varepsilon}
          +\frac{2L_{\rm hist}Y_*}{\varepsilon}\right).
\end{equation}
The lower bound on $a$ was used here. The constant $L_m$ bounds the
spatial Lipschitz norm of the normal at activation. The constant
$L_{\rm hist}$ bounds, in operator norm, the spatial derivative of the
linear map $Y_{\rm coord}\mapsto b_{Y_{\rm coord}}(t_0,y)$.
Finally $Y_*$ bounds the norm of the selected coordinate endpoint.
Consequently the spatial variation of its high velocity is at most
$\rho L_{\rm hist}Y_*$.

The reference ray is
\[
 y_0(\tau)=(\beta\tau^2,-2\beta\tau,1,0).
\]
Its full generator $C_0^*-2A_0^*=-p\otimes q-2\beta q\otimes n$
is nilpotent of order three. Its evolution is exactly
\[
 I-(\tau-s)p\otimes q-2\beta(\tau-s)q\otimes n+
                  \beta(\tau-s)^2p\otimes n
\]
and has norm at most $5\Theta^2$.
Duhamel and a supremum absorption now give, if
$60e\Theta^3\leq1$,
\begin{equation}\label{eq:ambient-ray-error}
 \|y-y_0\|\leq310e\Theta^5.
\end{equation}
Indeed the difference generator has norm at most $6e$; the resulting
bound is
$10\Theta^2(e+30e\Theta^3)\leq310e\Theta^5$.
Under the stronger smallness condition used below, this is at most
$1/2$. Thus
\[
 N\geq\tfrac12,\qquad |D_ry|^2\geq\tfrac14,\qquad
 \|y\|\leq6\Theta^2.
\]
In particular the extra ray components $r$ are controlled by
\eqref{eq:ambient-ray-error}; they have not been set to zero.

Tangency is $y\cdot z=0$, so the eliminated coordinate is exactly
\begin{equation}\label{eq:extra-tangent-elimination}
 W=-\frac{PU+QV+r\cdot Z}{N}.
\end{equation}
Let $\zeta=(U,V,Z)$, let $\Pi$ omit the $n$ coordinate, and let
$E_y\zeta$ insert the value \eqref{eq:extra-tangent-elimination}.
Then $\|E_y\|\leq1+2\|y\|\leq13\Theta^2$, and the actual reduced
equation on the full $d-1$ dimensional transverse space is
$\zeta'=\Pi\mathscr B_\varepsilon(y)E_y\zeta$.
It retains all couplings to $Z$.

The pressure denominator admits the following operator-norm estimate on
 the full ambient space. If
$r_*=\|y-y_0\|\leq1/2$, $R_*\geq\max(\|y\|,\|y_0\|)$ and
$\eta=\max(\|A-A_0\|,\|C-C_0\|)\leq1$, then
\begin{align}
 \|\mathscr B_\varepsilon(y)-\mathscr B_0(y_0)\|
 &\leq \eta+
 8R_*(4+8R_*^2)(r_*+\varepsilon R_*)\nonumber\\
 &\hspace{18mm}+2R_*(\eta R_*+3r_*).              \label{eq:ambient-generator-gap}
\end{align}
To prove it, put $u_\varepsilon=D_r^2y/|D_ry|^2$.
The preceding lower denominator bounds and
$\|D_ry-D_0y_0\|\leq r_*+\varepsilon R_*$
give
$\|u_\varepsilon-u_0\|\leq
(4+8R_*^2)(r_*+\varepsilon R_*)$.
Also
$\|A^*y-A_0^*y_0\|\leq\eta R_*+3r_*$,
$\|A^*y\|\leq4R_*$ and $\|u_0\|\leq R_*$.
Insert these estimates in the difference of the two rank-one
operators in \eqref{eq:ambient-exact-generator}.
All estimates are Euclidean operator estimates on the complete space.

Using \eqref{eq:ambient-ray-error}, $\eta\leq2e$,
$R_*\leq6\Theta^2$, and $\varepsilon\leq e$, the right side of
\eqref{eq:ambient-generator-gap} is at most
$5\cdot10^6e\Theta^{11}$.
The ideal generator has zero $n$ input column.
Therefore
$\Pi\mathscr B_0(y_0)(E_y-E_{y_0})=0$, and the actual reduced generator
differs from the ideal reduced generator by at most
\begin{equation}\label{eq:ambient-reduced-gap}
 65\cdot10^6e\Theta^{13}.
\end{equation}
The ideal reduced generator is explicitly
\[
 \begin{pmatrix}
 \dfrac{2P_0Q_0}{1+P_0^2}&
       -2+\dfrac{2P_0(P_0+\beta)}{1+P_0^2}\\[2mm]
 -1&0
 \end{pmatrix}
 \oplus0_{\mathcal E},
 \qquad P_0=\beta\tau^2,\quad Q_0=-2\beta\tau .
\]
Hence its active equation is
\[
 ((1+\beta^2\tau^4)V')'=2(1-\beta^2\tau^2)V,\qquad U=-V',
\]
and every ideal extra coordinate is constant.

Let $g=V_0$ be the positive scalar solution with $g(0)=1$, $g'(0)=0$.
The scalar relative propagator and the ratio estimate
$g(s)/g(t)\leq C\Theta$ bound the complete ideal reduced evolution by
\[
 K_d\frac{g(t)}{g(s)},\qquad K_d=C_d\Theta^8 .
\]
The extra identity block is included using this ratio estimate;
$g$ need not be nondecreasing. For completeness, whenever the reduced
coefficient gap is $\Delta$, weighted Duhamel gives
\[
 \frac{g(s)}{g(t)}\|U(t,s)\|
 \leq K_d+K_d\Delta
       \int_s^t\frac{g(s)}{g(r)}\|U(r,s)\|\,dr .
\]
Thus $K_d\Theta\Delta\leq1/2$ implies
$\|U(t,s)\|\leq2K_dg(t)/g(s)$.
This estimate covers arbitrary transverse initial data and forcing,
including additional modes.

It also bounds the true joined high evolution in the fixed transverse
coordinates used in the history problem. At each endpoint the actual
physical field and reduced coordinates are related by
\[
 b=O D_vE_y\zeta,\qquad
 \zeta=\Pi D_v^{-1}O^*b.
\]
Their operator norms are at most $13\Theta^2$ and
$\varepsilon^{-1}$, respectively. If $R_\perp:U\to m_0^\perp$
is an isometric identification and $b=F R_\perp a_{\rm hist}$, then
\[
 a_{\rm hist}=R_\perp^*F^{-1}b,\qquad
 b=F R_\perp a_{\rm hist}.
\]
These two endpoint changes cost at most $\|F^{-1}\|$ and $\|F\|$.
Thus the evolution in the actual $a_{\rm hist}$ coordinates is bounded
by $26K_d\Theta^2\varepsilon^{-1}
\sup\|F^{-1}\|\sup\|F\|\,g(t)/g(s)$.
Every factor is a known polynomial source cost. Applying the same
endpoint maps to the variation-of-constants formula gives the forced
tail bound with these costs. This is the full $d-1$ dimensional
forward operator joined to the history trace, including its action
on all additional transverse components.

For the selected primary let $\lambda\geq0$ be its actual normalized
initial slope and let $V_\lambda(0)=1$, $V_\lambda'(0)=\lambda$.
The ideal reduced initial vector is $(-\lambda,1,0)$.
At a neighboring label the discrepancy of $(U,V,Z)$ is at most
$\rho L_{\rm hist}Y_*/\varepsilon$, hence at most $e$ after the choice
\eqref{eq:ambient-error-input}. In particular the extra primary
initial component is bounded by the same history Lipschitz estimate. The remaining initial component is
determined by \eqref{eq:extra-tangent-elimination}.

Weighted Duhamel applied to the difference from
$\zeta_\lambda=(-V_\lambda',V_\lambda,0)$ now yields
\begin{equation}\label{eq:ambient-primary-relative}
 \|\zeta-\zeta_\lambda\|
 \leq C_de\Theta^{30}(1+\lambda)g(\tau).
\end{equation}
Indeed the initial term is $K_de\,g(\tau)$ and the integral term is
at most $2K_d^2\Theta\Delta(1+\lambda+e)g(\tau)$.
The exponent is $2\cdot8+1+13=30$.
Restoring the eliminated coordinate costs at most two more powers
of $\Theta$, by \eqref{eq:extra-tangent-elimination} and
\eqref{eq:ambient-ray-error}.

There is no loss of $\lambda$ on the late interval.
The second fundamental solution satisfies
\[
 \frac{V_1(\tau)-g(\tau)}{g(\tau)}
   =\int_0^\tau\frac{ds}{(1+\beta^2s^4)g(s)^2}.
\]
On $[0,1]$ the scalar equation has uniformly bounded coefficients,
so this integral at one has a fixed positive lower bound.
Consequently, for $\tau\geq1$,
$V_\lambda(\tau)\geq c(1+\lambda)g(\tau)$.
The actual normalized flux is $J=\langle y,Az\rangle$; the physical
flux for the unamplified field equals $s_0aJ$; the amplified field has
flux $\alpha s_0aJ$.
The ideal flux is
$J_0=Q_0(-V_\lambda')+(P_0+\beta)V_\lambda$.
Equations \eqref{eq:ambient-ray-error}--\eqref{eq:ambient-primary-relative}
give
\begin{equation}\label{eq:ambient-flux-relative}
 |J-J_0|/V_\lambda\leq C_de\Theta^{34}
 \quad(1\leq\tau\leq\tau_{\rm end}).
\end{equation}
For example, expand its difference into
$\langle y,(A-A_0)z\rangle+
 \langle y-y_0,A_0z\rangle+
 \langle y_0,A_0(z-z_\lambda)\rangle$.
The final term costs two powers of $\Theta$ beyond the restored-state
bound; the other two terms are smaller.

The scalar formula gives $J_0/V_\lambda\geq\beta=\sigma^2$.
To recall its margin, below $x=\sigma\tau=1$ both the
$Q_0(-V_\lambda'/V_\lambda)$ contribution and $P_0$ are nonnegative.
Above $x=1$ its equivalent expression is
$x^2-\beta+2\sqrt\beta\,z/x\geq1-\beta\geq3/4$, with the
nonnegative scalar Riccati variable $z$.
Let $x_{\rm tar}\geq2$ be the target coordinate and assume
$\tau_{\rm tar}=x_{\rm tar}/\sigma\leq\Theta$.
Then $\beta\geq4/\Theta^2$.
A single conservative guard
\begin{equation}\label{eq:ambient-sixty-guard}
 C_de\Theta^{60}\leq1
\end{equation}
therefore proves positive actual late flux, as well as all the ray,
denominator, reduced-propagator and relative-state guards above.
Enlarge the fixed $C_d$ to cover their finitely many numerical constants.

The same guard transfers the ideal uniform-size estimate to the actual
neighbor. Indeed the unamplified size is
$s_0\|D_ry\|\|D_vz\|$; the packet size is $\alpha$ times this quantity.
Subtract the corresponding central/reference expression for the packet size and use the
two state bounds. At the target the actual $q$ velocity component is at least
$V_\lambda(\tau_{\rm tar})/2$ and the normal component $N$ is at least
$1/2$. Hence the unamplified target size is at least
$s_0V_\lambda(\tau_{\rm tar})/4$.
The actual normalization $\alpha\,\mathrm{size}_{\rm target}
=\delta h_{\rm child}$ and the scalar ratio
$g(\tau)/g(\tau_{\rm tar})\leq C\Theta$ make this size error at most
$C_de\Theta^{40}\delta h_{\rm child}$, hence at most
$\delta h_{\rm child}$.
On $0\leq\tau\leq1$, the reference states are at most $C(1+\lambda)$
and the perturbation bound adds at most the same size.
The scalar target growth gives
$\alpha s_0(1+\lambda)\leq
 C\Theta\delta h_{\rm child}e^{-b/\sigma}$.
Here $b>0$ is a fixed constant from the scalar growth lemma, decreased
if needed. Thus the early size retains its exponentially small factor and a
fixed polynomial in $\Theta$ (the larger $\Theta^5$ bound suffices).
All these statements concern the actual short interval ending at
$\tau_{\rm end}$; the ideal scalar equation alone is defined globally.

\paragraph{The bounded profile height.}
The preceding normalization also closes the polynomial bound for the
profile constant $H_0$ through cancellation of the scalar growth
against the terminal normalization.
The forward scalar profile is $g=V_0$. In physical time this means
$\gamma(t)=\alpha$ for $t\leq t_0$ and
$\gamma(t)=\alpha V_0(a(t-t_0)/\varepsilon)$ for $t\geq t_0$.
Since $V_0\leq V_\lambda$,
the scalar ratio estimate on the retained interval through the target
and its stipulated short extension gives
$g(\tau)/V_\lambda(\tau_{\rm tar})\leq C\Theta$.
Together with the actual target-size lower bound this yields
\begin{equation}\label{eq:actual-profile-height}
 H_0=\max(1,\sup\gamma)
 \leq1+C\Theta\,\delta h_{\rm child}s_0^{-1}.
\end{equation}
The constant $s_0$ is that of the actual unit initial phase:
$s_0^{-1}=\|F(t_0)^*n_{\rm new}\|\leq\|F(t_0)\|$.
It is therefore an old particle-label cost.
Likewise the scalar target lower bound gives
\[
 \alpha\leq
 C\Theta\,\delta h_{\rm child}s_0^{-1}
              (1+\lambda)^{-1}e^{-b/\sigma}.
\]
Thus both the polynomial height and the small history amplitude come
from normalization at the same target; the resulting source majorant
has polynomial dependence on the parent parameters.

Finally, \eqref{eq:ambient-error-input} explains how to choose a
quantitative spatial neighborhood. Only old unnormalized coefficient
and history Lipschitz costs occur. Their positive-word inverse and
endpoint estimates are polynomial in the old particle-label data.
These costs depend only on the parent geometry, its selected history
and $\varepsilon^{-1}$. In particular, they are independent of the new
angular width $\delta_n$. The prescribed $\rho_n$ below satisfies
the resulting neighborhood conditions.

\subsection{One radius and polynomial source costs}

For an ordered-word Gevrey block put
$\mathfrak m_{R,b}(n)=R^{n+b}((n+b)!)^2$. The coefficient fixed-base
cost is bounded by the explicit polynomial
\[
 S_{d,q}(R_c,C)=2^q C\sum_{r=0}^q(4(d+1)R_c)^r(r!)^2.
\]
Here $R_c$ is the known coefficient radius. The full word alphabet is
$\{1,\ldots,d+1\}$.

Indeed, at external order $n$ and base order $0\leq r\leq q$,
\[
 (n+r)!^2\leq4^{n+r}(n!)^2(r!)^2.
\]
The $(d+1)^{n+r}$ ordered coordinate choices therefore give a
coefficient-block bound
\[
 S_{d,q}(R_c,C)\,[4(d+1)R_c]^n(n!)^2.
\]
Thus the known coefficient radius is replaced once by
$4(d+1)R_c$. In the commutator argument below $R_c$ denotes this
enlarged known radius. This conversion is applied only to coefficient
tensor bounds; the forcing is already measured by its full ordered-word
block.

For a positive continuous profile $\gamma$ and a fixed base order $q$,
the normalized block used here is
\[
 \mathcal B_n^\gamma(f)=
 \sum_{I\in\{1,\ldots,d+1\}^{n}}
 \bigl\|\gamma^{-1}\partial_I f\bigr\|_{C([0,T];H^q)}.
\]
The finite sum is outside the time norm. Replacing it by a supremum of
the sum is unnecessary in the induction. At each fixed time the block
dominates the corresponding sum of Sobolev norms and hence the
weighted norms used in the correction estimates. Equivalent choices
of the base $H^q$ norm change only constants depending on $d,q$.

\begin{lemma}[Same-radius joined bounds]\label{lem:joined-same-radius}
There is a threshold $R_*\geq1$, polynomial in the known source coefficients,
their reciprocal coercivity bounds, inverse time lengths and the relative
propagator constant of Section~\ref{subsec:full-transverse-relative}, such that
for every $R\geq R_*$ the joined high inverse maps a forcing satisfying
$\mathcal B_n^\gamma(f)\leq A\mathfrak m_{R,b}(n)$ to velocity and
time derivative bounds
\[
 \mathcal B_n^\gamma(h_f)\leq C_{\rm v}A\mathfrak m_{R,b+3}(n),\qquad
 \mathcal B_n^\gamma(\partial_t h_f)\leq C_{\rm t}A\mathfrak m_{R,b+3}(n).
\]
Here $h_f$ denotes the joined high field. The constants and radius
threshold are independent of $A$, $b$ and $R\geq R_*$. The
mean inverse, pressure inverse and tensor corrector have analogous
fixed-shift bounds at this same radius.
\end{lemma}
\begin{proof}
The history base estimate is Lemma~\ref{lem:joined-history-polynomial}.
Its terminal trace supplies the actual forward datum. The forward
base estimate is the relative estimate in Section~\ref{subsec:full-transverse-relative}, on the
fixed packet support. Products with the known frame turn coordinate
velocities into physical velocities, and the Gram equation controls
the true time derivative.

Here is the all-order induction, including the radius issue. Differentiating
the actual equation produces positive-order coefficient commutators.
The elementary inequality
\[
 \binom nj(j!)^2((n-j+b)!)^2\leq((n+b)!)^2
\]
holds for every $b\geq0$. If $M$ is the fixed-base inverse or forward
cost and $B R_c^j(j!)^2$ is the known coefficient block, the normalized
positive-order sum is bounded by
\[
 MB\sum_{j\geq1}(R_c/R)^j\leq\tfrac12
 \quad\hbox{when}\quad R\geq R_c(1+2MB).
\]
Thus a factor two closes the induction at all orders, without changing
the unknown forcing radius. The subsets of an ordered word give exactly
the binomial convolution above. A dimensional factor is permitted when
converting the \emph{known} coefficient tensor bounds into word bounds;
it is absorbed once into $R_c$.

The fixed-base inverse is the finite polynomial recursion in
Lemma~\ref{lem:joined-history-polynomial}. The mean form and pressure
equation have the same commutator identity and use their actual coercive
inverses. The corrector uses known coefficient products, one angular
primitive, and one spatial derivative. The primitive costs at most
the fixed period and the coordinate sums are finite ($\leq d^3$
terms). Taking the sum of these finitely many required radii gives
a polynomial common radius. Linearity, applied after dividing the
forcing and datum by $A$, proves that the radius is independent of
$A$; the zero-amplitude case is uniqueness. The time profile itself
is never differentiated.
\end{proof}

These estimates define an explicit polynomial source majorant. Replace
known primitive bounds by $W\geq1$, coefficient blocks by $S_{d,q}$,
fixed inverses by their $q$-step polynomial recursion, and maxima by
sums. Reciprocal coercivity, ray and Gram constants are first replaced
by their \emph{proved} bounds. Every remaining operation in the
initialized source construction is addition, multiplication, a fixed
power or one of these fixed recursions. The resulting positive
polynomial $\Phi_d(W)$ gives the bound
\[
 \Phi_d(W)\leq A_dW^{e_d},\qquad
 e_d=\deg\Phi_d,\qquad A_d=1+\sum_j|[W^j]\Phi_d|.
\]
The exponent is therefore determined by a fixed finite computation.
The mean and high grade recursion uses this common radius at every grade.

\subsection{The explicit nonlinear threshold}

Let $p\geq300$ be a fixed tail-polynomial degree and
$\vartheta=1/(1000(p+1))$. The grade bounds in the
finite recursion permit $p=300$. Write
$Z=k^\vartheta$, $N=\lfloor Z\rfloor$. Let $C,I,G,D,T,r_{\mathrm a}^{-1}$
be the actual tail, inverse, correction-growth, drift, time and reciprocal
radius costs of that packet, with $C,I,G,D,T\geq0$.
Here $r_{\mathrm a}>0$ is its analytic
radius; it is distinct from the physical localization length $\rho$ and
from the fixed mean cutoff parameter $\ell_*$. Define
\begin{equation}\label{eq:source-cap}
 H=64+\exp(1)+2^pC+I+12GT+\frac{8GTD}{r_{\mathrm a}}
                   +\left(\frac{8GT}{r_{\mathrm a}}\right)^2,\qquad
 k_*=H^{1000(p+1)}.
\end{equation}
For $k\geq k_*$ one has $H\leq Z$, $Z\geq64$ and $\log k\geq1$.
In particular
\[
 C(N+1)^p\leq C(2Z)^p\leq Z^{p+1}=k^{1/1000}.
\]
The genuine residual and correction tolerance obey
\[
 2[2e^{-(7/10)Z\log k}]e^{3GT}\leq\tfrac12e^{-\sqrt Z},
 \qquad
 2G(D/k+e^{-\sqrt Z})T\leq r_{\mathrm a}/2.
\]
Indeed $3GT\leq Z/4$, $\sqrt Z\leq Z/8$, $\log8\leq Z/8$ prove
the first inequality. For the second, use
$8GTD\leq r_{\mathrm a}k$ and
$8GT/r_{\mathrm a}\leq\sqrt Z\leq e^{\sqrt Z}$.
This gives a polynomial frequency threshold for the exponential
correction-growth condition.

If the six costs are bounded by $Q\geq1$, the explicit bound is
\begin{equation}\label{eq:polynomial-frequency-threshold}
 k_*\leq(2^p+200)^{1000(p+1)}Q^{6000(p+1)}.
\end{equation}
For $p=300$ the power of $Q$ is $1\,806\,000$. Its size is irrelevant;
its independence of the stage and the truncation is essential.

\subsection{The actual next particle-label bound}

Here $u$ denotes the coordinate packet in normalized parent labels, and
$m_0$ is the constant initial covector. The corrected packet's lifted velocity is
\[
 \mathcal V=(k^{-1}u,\langle m_0,u\rangle),\qquad \|m_0\|=1.
\]
Its normal cancellation and the drift energy bound give supremum and
cylinder $L^2$ jet amplitudes
\[
 B\leq C_d(D/k+C_{\rm e}e^{-\sqrt{k^\vartheta}}),
 \qquad R_*\leq C_d/r_{\mathrm a}.
\]
The actual time derivative has amplitude $C_1$, a polynomial source
cost, at a fixed enlargement of $R_*$. Choose a fixed
$\chi<\min(1/100,\vartheta/100)$. The source comparison below yields
$C_dD,C_dC_{\rm e},R_*,T,C_1\leq k^\chi$. A fixed numerical lower
bound on $k$ then gives
\[
 B\leq2k^{-1/2},\qquad BR_*T\leq1/8.
\]
For example, putting $a=\vartheta/2$ and $v=\log k$, the inequality
$e^{av}\geq a^2v^2/2$ shows that
$v\geq2(2+\chi)/a^2$ implies
$k^\chi e^{-\sqrt{k^\vartheta}}\leq k^{-2}$.

Lemmas~\ref{lem:graph-invariant} and~\ref{lem:small-lift}, the
one-dimensional phase trace, and the affine graph derivative
factor $1+k$ now give the actual relative-label displacement,
velocity and acceleration $(a,b,c)$:
\[
 \|D^r a\|_{2,\infty},\ \|D^r b\|_{2,\infty}
 \leq k^{-1/4}(\rho^{-1}k^{5/4})^r(r!)^2,\qquad
 \|D^r c\|_2\leq k^{1/4}(\rho^{-1}k^{5/4})^r(r!)^2.
\]
The time acceleration is the actual field
$(\mathcal V_t+D\mathcal V\,\mathcal V)\circ\Psi$.
The only small-time flow condition is $BR_*T\leq1/8$ on the
lifted field.

To record the polynomial degree, let $X=\mathrm{id}+D_0$ be the
old particle map, with actual velocity $V_0$ and acceleration $W_0$,
and $Y=\mathrm{id}+a$ the relative map. Then
\[
 \begin{split}
 D_0^+&=D_0\circ Y+a,\\
 V_0^+&=V_0\circ Y+b+(DD_0\circ Y)b,\\
 W_0^+&=W_0\circ Y+2(DV_0\circ Y)b+
        (D^2D_0\circ Y)[b,b]+c+(DD_0\circ Y)c.
 \end{split}
\]
Suppose the old particle fields have common Gevrey amplitude and radius
at most $K\geq1$, with displacement and velocity controlled in both
$L^2$ and supremum norm and acceleration controlled in $L^2$.
A sufficient new amplitude
and radius, with $A=k^{1/4}$ and $R=\rho^{-1}k^{5/4}$, are
\[
 A^+=K+A+9C_dK^2A+36C_dK^3A^2,\qquad
 R^+=(1+R)((1+A)16K+2)+R .
\]
These follow term by term from the factorial product and composition
inequalities. Volume preservation permits composition in $L^2$.
If $K,C_d,\rho^{-1}\leq k^{1/100}$, fixed numerical margins give
$A^+\leq k$, $R^+\leq k^2$. A fixed Sobolev shift gives the next
particle-label bound $K^+\leq k^{2q+4}$. In particular the acceleration
has supremum jet bounds after that shift, by the fixed-dimensional
Sobolev embedding. This supplies the coefficient bounds for $F_2^+$
used in the next source.

The full inverse deformation is polynomial as well. Since
$\det F^+=1$, its inverse is the adjugate. If
$\|D^nF^+\|\leq A R^n(n!)^2$, expansion of each cofactor into
$(d-1)!$ products and the Gevrey product constant 3 give
\[
 \|D^n(F^{+,-1})\|
 \leq d(d-1)!\,3^{d-2}A^{d-1}R^n(n!)^2.
\]
The factor $d$ bounds operator norm by the largest entry.
The time derivative satisfies
$(F^{+,-1})_t=-F^{+,-1}F_1^+F^{+,-1}$.
Also $F_1^+=D_xV_0^+$ and $F_2^+=D_xW_0^+$.
Thus the source coefficients have fixed polynomial bounds in the
new particle-label parameter. Fix, for instance,
\[
 D_d=100(q+d+1)
\]
for the particle-label invariant $K^+\leq k^{D_d}$.
The inverse and coefficient costs are fixed polynomials in $K^+$;
they are not asserted to have that same exponent $D_d$.

\subsection{Uniform scales, with no second frequency choice}

Set $s_n=J+n$ and
\[
 x_0=X,\quad x_{n+1}=s_n^2x_n,\quad
 k_n=e^{x_n/s_n^2},\quad h_n=e^{x_n/s_n^5},\quad
 \delta_n=e^{-x_n/s_n^3},\quad\rho_n=e^{-x_n/s_n^{7/2}}.
\]
Here $n=j-J$ and $x_n=x_{j-1}$ in the notation of
Section~\ref{sec:iteration}; $\rho_n$ is the physical localization,
not the fixed mean cutoff $\ell_*$. For $n\geq1$ the identities
\[
 \log k_{n-1}=\frac{x_n}{(s_n-1)^4},\quad
 \log h_{n-1}=\frac{x_n}{(s_n-1)^7},\quad
 \log\rho_{n-1}^{-1}=\frac{x_n}{(s_n-1)^{11/2}}
\]
are exact. If $K_n\leq k_{n-1}^{D_d}$, every old geometric cost in
\eqref{eq:ambient-error-input} is bounded by
$Cs_n^a x_n^b e^{c x_n/(s_n-1)^4}$.
For sufficiently large fixed $J$ and then $X$,
\[
 \rho_n Cs_n^a x_n^b e^{c x_n/(s_n-1)^4}
 \leq e^{-x_n/(2s_n^{7/2})}.
\]
The additional factors in the ambient guard have the same inherited
scale: $\Theta_n\leq C(1+s_n^2x_n^2)$,
$s_0^{-1}\leq\|F(t_0,0)\|$, and
$\varepsilon^{-1}\leq C\sqrt{h_{n-1}}$.
They therefore contribute only polynomial factors in $s_n,x_n$ or an
exponential with denominator $(s_n-1)^7$, which is already absorbed
by the displayed denominator $(s_n-1)^4$. The terms
$\varepsilon\Theta G^2$ and $e_{\rm old}$ in
\eqref{eq:ambient-error-input} are controlled by the separate shear
separation and old frequency-error comparisons in
Section~\ref{sec:iteration}; they are not made small by localization.
At $n=0$ the activation time is zero, and the forward high construction
is used. This branch contains no inverse-history-time cost.
The seed predecessor has polynomial costs in $X$. The normal-stage factor $e^{-X/J^{7/2}}$ dominates those costs for the
same final large choice of $X$.

Consequently all neighborhood conditions hold at the prescribed
localization length $\rho_n$.

The full source list additionally includes $\delta_n^{-1}$ and the
new scalar height. Its positive-polynomial bounds therefore have the
form
\begin{equation}\label{eq:uniform-parameter-envelope}
 \mathcal P_n=Cs_n^a x_n^b
              \exp\left(\frac{c x_n}{(s_n-1)^3}\right).
\end{equation}
In particular the radius, $H_0$, reciprocal analytic radii, history
and mean inverse constants, and the costs in \eqref{eq:source-cap}
are bounded by fixed powers of \eqref{eq:uniform-parameter-envelope}.

\begin{lemma}[Uniform domination]\label{lem:uniform-source-domination}
For a finite list of fixed $A,C,c,a,b,N>0$ and $\gamma>0$, one common
choice of $J$ and then $X$ ensures
\[
 A\mathcal P_n^N\leq k_n^\gamma\qquad(n\geq0)
\]
for every entry of the list.
\end{lemma}
\begin{proof}
Choose $J$ so that
$Nc/(s-1)^3\leq\gamma/(2s^2)$ for $s\geq J$;
$J\geq16Nc/\gamma+2$ suffices. Increase $X$ until
\[
 \log(AC^N)+Na\log s_n+Nb\log x_n
 \leq\frac{\gamma x_n}{2s_n^2}
\]
uniformly in $n$. To see the uniformity, use
$x_n=X\prod_{r<n}(J+r)^2$.
The sequences
$s_n^2\log s_n/\prod_{r<n}(J+r)^2$ and
$s_n^2\log(\prod_{r<n}(J+r)^2)/\prod_{r<n}(J+r)^2$
are bounded, since the denominator eventually dominates every
polynomial in $n$ while its logarithm grows at most
$2n\log(J+n)$. The resulting bound is
$O_J((1+\log X)/X)$ and tends to zero.
Adding the two inequalities proves the assertion. Take the maximum
of the finitely many choices of $J,X$.
\end{proof}

Apply this lemma simultaneously to the source costs, the neighborhood
guards, \eqref{eq:polynomial-frequency-threshold}, the small lifted-flow
costs and their numerical margins. It gives the actual nonlinear
correction at the preselected $k_n$ and then
$K_{n+1}\leq k_n^{D_d}$. The exponent $D_d$ is fixed first;
the source polynomials and their degrees are then fixed; $J,X$ are
chosen once for the resulting finite list. No exponent or starting
index is enlarged after a later child is observed.

The induction applies to the joined variational-history family and
the quantitative neighborhood constructed in this section.
